\documentclass[11pt]{article}
\usepackage[T1]{fontenc}
\usepackage[margin=1.05in]{geometry}
\usepackage{amsmath,amssymb,amsthm,microtype}
\usepackage[hidelinks]{hyperref}
\newtheorem{theorem}{Theorem}
\newtheorem{proposition}[theorem]{Proposition}
\newcommand{\F}{\mathbb F_2}
\newcommand{\tr}{\operatorname{tr}}
\title{The classical bias of hyperbolic Pauli line constraints}
\author{}
\date{Research audit note, 25 November 2026}
\begin{document}
\maketitle
\begin{abstract}
For the real $n$-qubit Pauli quadric, let $D$ be the minimum number of violated lines and $L$ their total number. With $d=2^n$, the difference $1/2-D/L$ is at most $\sqrt{2(d+2)}/(d+4)$. The cubic spectral floor of Jonah gives this bound directly. We give an elementary proof, a second derivation from his section theorem, and an explicit randomized assignment giving a lower bound of order $d^{-1/2}$ on the same difference. Thus the rate of approach to $1/2$ is sharp in order. The original upper bound on the gap and its limit have an end-to-end Lean proof. We distinguish these results from earlier four- and five-qubit classifications.
\end{abstract}

\section{The limiting question and its prior answer}
The revised manuscript of Li, Gupta and Ganesh asks in its fourth open problem for the limit of the violated fraction on higher-rank hyperbolic quadrics \cite{lgg,lggrepo}. Its version on arXiv on 19 October 2026 asks a different higher-rank question. An October preprint of Jonah already bounds the signed excess of satisfied lines by the solution of a cubic moment problem \cite[Thms.~1,6,8]{structure}. We derive the requested limit as a corollary and determine the order of its error. The separate equality question for adjacent hyperbolic and elliptic sections remains.

Let $E=\F^n$, $V=E\oplus E$ and $q(a,b)=a\cdot b$. A point is a nonzero $p=(a,b)$ with $q(p)=0$. A line is an unordered triple $\{p,r,p+r\}$ of points. Write
\[
 P(a,b)e_x=(-1)^{b\cdot x}e_{x+a}.
\]
Write $E_p=P(p)$. Since $E_p^T=(-1)^{q(p)}E_p$, the quadric labels the traceless symmetric Paulis. For $r=(c,h)$ we have $E_pE_r=(-1)^{b\cdot c}E_{p+r}$. On a line these operators commute, since $a\cdot h+b\cdot c=q(p+r)+q(p)+q(r)=0$. Their product is $\epsilon_\ell I$. Individual changes of Pauli representatives multiply the corresponding assigned signs by the same factors and preserve the violated set. The usual tensor convention uses $Y=iXZ$ and changes the representative by $(-1)^{k/2}$ when $k$ factors carry both an $X$ and a $Z$.

An assignment $s$ chooses one sign at each point. A line is violated when $\epsilon_\ell\prod_{p\in\ell}s_p=-1$. Set $D=\min_s N(s)$, where $N(s)$ counts such lines. The maximum signed bias is $\beta=1-2D/L$.
\begin{theorem}\label{main}
For every $n\ge2$,
\[
 0\le \frac12-\frac DL\le \frac{\sqrt{2(d+2)}}{d+4}.
\]
Moreover,
\[
 \sqrt2 e^{-3/2}\le
 \liminf_{n\to\infty} \sqrt d\left(\frac12-\frac DL\right)
 \le\limsup_{n\to\infty} \sqrt d\left(\frac12-\frac DL\right)\le \sqrt2.
\]
\end{theorem}

\section{The spectral bound}
Character orthogonality gives $\sum_{a,b}(-1)^{a\cdot b}=d$. Thus the number of points is $m=(d-1)(d+2)/2$. For a fixed point $p$, sum
$(1+(-1)^{q(r)})(1+(-1)^{q(r+p)})$ over $r\in V$. The constant contributes $d^2$, the two single characters contribute $2d$, and the product sums to zero by nondegeneracy of the polar form. Dividing by four and removing $r=0,p$ gives $(d-2)(d+4)/4$ partners. Each line has six ordered pairs, so
\[
 L=(d-1)(d-2)(d+2)(d+4)/48.
\]
Define $A=\sum_{p\in Q}s_pE_p$. Its trace vanishes, and $\tr A^2=dm$ by Pauli orthogonality. A triple in $\tr A^3$ has zero trace unless its labels sum to zero. The three nonzero singular labels then constitute a line. All six permutations contribute the same sign. Hence
\begin{equation}\label{trace}
 \tr A^3=6d(L-2N(s)).
\end{equation}
This identity also follows from Jonah's operator identity for sets closed under commuting products \cite[Thm.~1]{structure}.

Put $a=\sqrt{(d+2)/2}$ and let $\lambda_i$ be the eigenvalues. Their sum is zero and $\sum\lambda_i^2=d(d-1)a^2$. Cauchy--Schwarz on the coordinates other than $i$ gives $\lambda_i^2\le(d-1)(d(d-1)a^2-\lambda_i^2)$, and thus $\lambda_i\le(d-1)a$. Expanding the sum of the nonnegative terms $((d-1)a-\lambda_i)(\lambda_i+a)^2$ gives
\begin{equation}\label{floor}
 \tr A^3\le d(d-1)(d-2)a^3 .
\end{equation}
This is the trace-zero cubic inequality of Okumura \cite{okumura}. Substitution in \eqref{trace} gives the first upper bound in Theorem~\ref{main}. Negating all assigned signs exchanges satisfied and violated lines, giving $D\le L/2$. The explicit error tends to zero and gives the limsup.

Jonah's explicit formula is $\mathrm{Spec}_Q=m^{3/2}(d-2)/(6\sqrt{d-1})$ with $L-2D\le\mathrm{Spec}_Q$ \cite[Thm.~6]{structure}. Algebraically it is precisely \eqref{floor}. His earlier section theorem gives another derivation \cite[Thm.~2]{section}. Set $h=d/2$. For an off-point hyperplane the section is a full $(n-1)$-qubit space. The $h(h+1)$ outside points support $L_O=h(h^2-1)(h+2)/4$ outside lines, by subtracting the section's $(h^2-1)(h^2-4)/12$ lines. For fixed section signs, the outside signed excess is $X^T\Lambda X/2$, where the largest eigenvalue of $\Lambda$ equals that of a signed full Pauli sum $A'$ in dimension $h$. Since $\tr A'=0$ and $\tr((A')^2)=h(h^2-1)$, its largest eigenvalue is at most $(h-1)\sqrt{h+1}$. Division by $L_O$ bounds the outside bias by $2\sqrt{h+1}/(h+2)$. The full orthogonal group is transitive on lines and permutes these hyperplanes. Every line occurs equally often outside, so averaging gives the same global bound.

\section{An assignment attaining the order of the error}
Fix $t>0$ and draw a standard Gaussian vector $G$. Given $G$, choose the signs independently, with conditional means
\[
  \mathbb E[s_p\mid G]=\sin\!\left(t G^TE_pG/\sqrt d\right).
\] The optimum is at least the expected signed bias. Fix a line and denote its first two operators by $P,Q$. Oddness of sine absorbs the line sign into the third operator, giving $R=PQ$.

The four common eigenspaces for $P,Q$ have dimension $d/4$: their projectors are $(I\pm P)(I\pm Q)/4$ and $P,Q,PQ$ have trace zero. Write $U,V,W$ for the three quadratic forms divided by $\sqrt d$. Among the four combinations entering the sine product identity, $U+V+W$ has coefficients $3,-1,-1,-1$ on these spaces. Each of $U+V-W,U-V+W,-U+V+W$ has the opposite multiset. Rotational invariance gives
\[
 \mathbb E[\sin(tU)\sin(tV)\sin(tW)]
 =-\operatorname{Im}\left[
 (1-6it/\sqrt d)^{-d/8}(1+2it/\sqrt d)^{-3d/8}
 \right].
\]
This uses the characteristic function $(1-2iz)^{-1/2}$ of a squared standard Gaussian. The right side is common to all lines.

For fixed $t$, the logarithm of the bracket is $-3t^2-8it^3/\sqrt d+O_t(d^{-1})$. The real part of the remainder is $O_t(d^{-1})$ and its imaginary part is $O_t(d^{-3/2})$, by Taylor expansion of the two logarithms near one. Consequently
\[
 \beta\ \ge\ 8t^3e^{-3t^2}/\sqrt d+O_t(d^{-3/2}).
\]
Taking $t=1/\sqrt2$ and dividing by two gives the lower liminf in Theorem~\ref{main}. No assumption on convergence of the rescaled gap is needed.

\section{Equality, physical scope and verification}
Equality in \eqref{floor} requires one eigenvalue $(d-1)a$ and $d-1$ copies of $-a$. Since $L-2D$ is integral, equality with $d\ge4$ forces $a$ rational and hence integral. Then $(a-1)(a+1)=2^{n-1}$. The two positive powers of two differing by two must be 2 and 4. Thus $n=4$, as also follows from Jonah's more general classification \cite[Thm.~8]{structure}. At $d=16$ the vector $(3,-1,\dots,-1)/\sqrt{24}$ has symmetric Pauli expectations of modulus $1/3$. Its signs give $A=48\psi\psi^T-3I$, hence $D=315$. The 61,440 assignments and 240 configurations of type $DW(5,2)$ were established by Li, Gupta and Ganesh following the heuristic discovery of Muller et al.\ \cite{lgg,muller}. Jonah had announced the flat-state equivalence and a conditional geometric argument before this project \cite{digest}. The principal-minor map from three-qubit Lagrangians is also established \cite{plucker}. A direct coordinate proof of the orbit, counts and incidence correspondence accompanies this audit.

The ideal, state-independent signed correlator averaged over all lines is one for quantum measurements and at most $\beta$ for deterministic noncontextual assignments and their mixtures. These statements concern the full compatible context family with uniform weights. A single-neutron experiment with spin, path and energy has realized a different, three-degree-of-freedom Mermin witness \cite{neutron}. Experimental use of the present family would require calibrated joint measurements and a treatment of compatibility and sampling. The theorem alone gives no experimental detection threshold for an implemented instrument.

The accompanying Lean 4.19.0 project defines the actual Pauli matrices, unordered lines and finite minimum. It checks the first inequality of Theorem~\ref{main}, its limit, and the four-qubit degree. The Gaussian construction is a mathematical proof outside that formal verification. The exact rescaled constant, possible oscillations, and explicit deterministic assignments attaining an optimal leading constant remain to be determined. Priority for the matching-order construction and for the detailed coordinate classification remains unsettled.

\bibliographystyle{plain}
\bibliography{references}
\end{document}
